Four measurements explain why the earlier replication attempt on this codebase (best return 9.20 and 378-step episodes after 10M steps on flat ground, versus $\approx$86 and 4000 steps reported by Salehi \cite{salehi2025}) did not converge. They were obtained with stock MuJoCo 3.13, the configuration used for all experiments below.

\textbf{(1) Wheel friction frames.} Table~\ref{tab:frame} lists the alignment between each wheel's capsule axis and the first contact tangent $t_1$; only wheel~0 is correct on stock MuJoCo.

\begin{table}[h]
\centering\caption{$|\langle \text{capsule axis},t_1\rangle|$ for the ball--wheel contacts (1 = correctly aligned).}\label{tab:frame}
\small\begin{tabular}{lccc}\toprule
 & Wheel 0 & Wheel 1 & Wheel 2\\ \midrule
Stock MuJoCo 3.13 & 0.999 & 0.209 & 0.537\\
With our frame fix & 1.000 & 1.000 & 1.000\\ \bottomrule
\end{tabular}\end{table}

\textbf{(2) Velocity indexing.} Comparing MuJoCo's \texttt{cvel} with finite differences of the base position, the last three entries are the linear velocity (e.g.\ $(0.017,-0.005)$ vs.\ $(0.016,-0.005)$\,m/s) while the first three are the angular velocity. The environment used the first three as ``velocity''.

\textbf{(3) Terrain steepness.} Table~\ref{tab:slope} gives the slope statistics of the height field (20 held-out seeds per row). The default amplitude of 1 produces slopes that no policy with a $20^\circ$ tilt limit can climb; the amplitudes used in our experiments span $1^\circ$--$12^\circ$ mean slope.

\begin{table}[h]
\centering\caption{Perlin terrain slope versus amplitude $A$ (degrees).}\label{tab:slope}
\small\begin{tabular}{rccc}\toprule
$A$ & mean slope & 95th pct. & max\\ \midrule
0.02 & 1.3 & 2.3 & 3.7\\
0.04 & 2.7 & 4.7 & 7.3\\
0.08 & 5.3 & 9.3 & 14.4\\
0.12 & 7.9 & 13.8 & 21.1\\
0.14 & 9.2 & 15.9 & 24.2\\
0.18 & 11.8 & 20.2 & 30.0\\
1.00 (default) & 46.2 & 63.9 & 72.7\\ \bottomrule
\end{tabular}\end{table}

\textbf{(4) Effect of the fixes on flat ground.} Table~\ref{tab:abl} shows what each repair changes on flat ground with identical PPO settings (proprioception only, 0.3M decisions, deterministic evaluation episode of 8\,s). The released-code row is the best prior replication (10M physics steps at 500\,Hz). Removing the friction fix leaves the robot unable to balance (it falls after 0.36\,s), whereas with the fix the policy balances for the whole episode. The velocity bug did not prevent balancing but removed any drive to move: the policy stood still (net displacement $-0.016$\,m), while the corrected reward gave $1.24$\,m in the same time and budget (an earlier, otherwise identical configuration without the fall penalty).

\begin{table*}[t]
\centering\caption{Effect of the repairs on flat ground (single seed).}\label{tab:abl}
\scriptsize\begin{tabular}{lccc}\toprule
Setup & Steps survived (of 800) & Fell & Net progress [m]\\ \midrule
Prior replication (released code)$^\dagger$ & 378 of 4000 (train mean) & yes & --\\
Stock friction frames, velocity fixed & 36 & yes & 0.00\\
Friction fix, velocity bug (swapped \texttt{cvel}) & 800 & no & $-0.02$\\
Friction fix, velocity fixed (all repairs) & 800 & no & 0.27--1.24$^\ddagger$\\ \bottomrule
\end{tabular}\\[2pt]
{\scriptsize $^\dagger$ training-time mean episode length at 500\,Hz. $^\ddagger$ 1.24\,m without the fall penalty and $\gamma{=}0.995$ at 0.3M decisions; 0.27\,m with the final setting ($\gamma{=}0.999$, fall penalty 10) at the same budget, which trades early speed for stability.}
\end{table*}

